\section*{Introduction}

% Owner: papers/1_method/decisions.md ("Introduction, five moves", 2026-07-28) and
% papers/_program/program.md §1 (merged macro paper, least squares as the comparison anchor).
% Written after 04_results.tex so that the ground conceded here is the ground the paper delivers.
% The concessions in the prior-art paragraph are sourced in
% docs/bibliography/temporal_correlation_and_AR_errors_2026-07-28.md, Part 2.

A macroscopic ion-channel current is the summed signal of a population of channels. Its mean carries the gating kinetics through the deterministic relaxation, and its fluctuations carry, in addition, the number of channels, the unitary current, and, through the temporal structure of the noise, the rate constants themselves. The standard way to fit a kinetic scheme to such a recording is least squares on the mean current. The four approaches surveyed by \citet{clerx2019four} are all deterministic fits with a root-mean-square objective, and the community benchmark IonBench states that it does not cover optimisation approaches for stochastic models \citep{owen2025ionbench}.

Least squares rests on an assumption that a gating record does not satisfy. The residual sum of squares becomes a variance, and then a confidence interval, only through a count of independent residuals. Channel gating leaves the residuals correlated over the channel's own relaxation time, so a sample taken faster than that time partly repeats what the sample before it already said. The point estimate often survives this. The count does not: a fit that treats correlated residuals as independent behaves as though the recording supplied more observations than it did, and reports an interval narrower than the spread of the estimator it is describing. \citet{delcore2025parameter} state the same cost in the field's own terms, that standard methods built for deterministic models do not distinguish stochastic channel gating from measurement noise and return biased estimates.

A Markov description of gating models the dependence instead of assuming it away. A channel with a small number of conformational states, connected by rate constants that hold over the whole record, generates the mean current and the covariance of the fluctuations from the same constants. The temporal structure of the noise then becomes a second measurement of the same parameters rather than a nuisance to be neutralised. The rate constants also mean something outside the record: they are properties of the protein, comparable across conditions, preparations and laboratories, and they connect what is measured to a mechanism. A correlation function fitted freely to the residuals delivers neither of those.

Using the fluctuations to identify a scheme requires a likelihood for the recorded current, and the likelihoods that exist for macroscopic recordings form a small lineage, each of them an approximation. \citet{celentano2004use} fit the full Q-matrix covariance directly, at a cost that grows as the cube of the record length. \citet{milescu2005maximum} keeps the gating variance at each sample and drops the correlation between samples. \citet{moffatt2007estimation} recasts the problem as a recursive filter, propagating the occupancy distribution and conditioning it on each new datum. \citet{stepanyuk2011efficient} restores the full covariance at near-linear cost through a semiseparable representation, and \citet{munch2022bayesian} places the recursive filter in a Bayesian setting. A recent member conditions the interval-averaged observable on the channel states at both ends of the acquisition interval (MacroIR, \citealt{moffatt2025bayesian}). Ordered by what they compute for every interval of the record, these methods form a ladder of cost with least squares at its bottom rung, and the ladder hangs from one root question: does the analysis model the gating fluctuations at all, and if it does, how much of their temporal structure does it keep?

None of these methods has displaced least squares. They are published, implemented and available, and they are almost unused. The reason has to do with what an analyst can see.

Least squares is checkable by eye. It predicts the current at every point of the record from the parameters alone, so the prediction can be laid over the data and read interval by interval. A systematic departure of the model from the record shows up as a systematic departure of the line from the points. When the scheme is wrong the picture usually says so, and when the forward model has a mistake in it the picture usually says that too. This is weak evidence in a formal sense and strong evidence in practice, and it is what has kept the classical fit in place.

A recursive filter offers no such picture. It predicts each interval from the parameters and from the data already observed, so the previous sample enters the prediction of the next one and the predicted trace follows the record closely by construction. It follows the record closely when the kinetic scheme is right, when the kinetic scheme is wrong, and when the code has a mistake in it. The residuals come out small and close to white, and they come out white in large part because the recursion subtracts, at each step, what it has just seen. Whiteness of the residuals is as far as the prior work goes, and it cannot falsify anything, since enough free terms will whiten any sequence. The visual evidence that carries the classical fit is unavailable for a filter, and nothing has taken its place. There has been no straightforward way to tell a working method from a bug, and no criterion for deciding whether the confidence interval it reports, which is the reason to use a likelihood in the first place, should be believed. Supplying that criterion is what this paper is for.

The standard statistical repair for correlated residuals does not reach this problem. Generalized least squares with autoregressive or moving-average errors whitens the residual sequence with a correlation function estimated from the residuals themselves, and \citet{lei2020considering} have run exactly that on ion-channel data, fitting an autoregressive moving-average error model and two Gaussian-process variants to the residuals of a hERG kinetic model. Gating puts this family out of reach for structural reasons. An autoregressive moving-average error model is stationary by construction, one covariance function for the whole trace, while the gating covariance tracks the mean current, vanishes when the current does, and restarts at every concentration or voltage jump. It contains no channel number and no unitary current, so it cannot return the two parameters the fluctuations are richest in. Its timescales are free parameters, where the Markov model ties them to the eigenvalues of the rate matrix that already fixed the mean, so free timescales absorb the information the mechanism would have contributed. \citet{lei2020considering} report that the autoregressive model widens the posterior without predicting better, and that in some cases the best predictions came from ignoring the discrepancy altogether. Their target is structural error in the kinetic scheme rather than finite-channel gating noise, so their study is not a test of the present question; what it shows is what a correlation model without a mechanism can and cannot buy.

% PENDING-BIB: the first prior-art citation of the paragraph below. katz1970membrane,
% katz1972statistical, anderson1973voltage, conti1980conductance, sigworth1981covariance and
% lei2020considering (cited above) live in biblio_full.bib (124 entries) and are absent from
% biblio.bib (15 entries), which is what elife_paper.tex currently points at. Repoint
% \bibliography{biblio} to biblio_full, or merge the six keys across. Not a retrieval job.
None of this is a claim that the correlation went unnoticed. Kinetic information has been read out of the temporal correlation of macroscopic currents since the early 1970s, first through the Lorentzian power spectrum under stationary conditions \citep{katz1970membrane,katz1972statistical,anderson1973voltage}, which is the Fourier transform of the autocovariance. \citet{anderson1973voltage} obtained the unitary conductance and the channel closing rate jointly, from the same fluctuation data and with standard errors on both: a unitary conductance of $0.205 \pm 0.0063$ (standard error) $\times 10^{-10}$ mho, and a closing rate exponential in membrane potential whose prefactor is $0.17 \pm 0.04$ ms$^{-1}$. % src: docs/bibliography/temporal_correlation_and_AR_errors_2026-07-28.md §Part 2 Tier 1, from Anderson & Stevens 1973 J. Physiol. 235:655-691 (PDF in docs/bibliography/)
The exact non-stationary two-time covariance for a population of identical independent Markov channels was derived and used a decade later \citep{conti1980conductance,sigworth1981covariance}, estimated from ensembles of $256$ to $504$ repeated sweeps, from which rate constants were obtained with standard errors, and used principally to discriminate between kinetic hypotheses. % src: docs/bibliography/temporal_correlation_and_AR_errors_2026-07-28.md §Part 2 Tier 2, from Sigworth 1981 Biophys. J. 34:111-133 (PDF in docs/bibliography/)
A covariance likelihood for macroscopic currents \citep{celentano2004use} predates the recursive filter of \citet{moffatt2007estimation} by three years. The statistical apparatus used below is equally old: the score, the information matrix equality and the sandwich covariance under misspecification are classical \citep{godambe1960optimum,huber1967behavior,white1982maximum}. What none of this work supplies is a statement of when a method built on the correlation is working.

Non-stationary fluctuation analysis is the one fluctuation-based method in routine use. It regresses the variance of the current against its mean and returns the unitary current, the channel number and the peak open probability \citep{stepanyuk2014maximum}, and practitioners of the method report that the unitary current is close to the only parameter it recovers reliably and that kinetic rates have essentially never been estimated this way for receptors in their native setting. This paper does not compete with it over how much amplitude information a macroscopic recording holds. The question here is whether a reported confidence interval from a fit to such a recording is honest, and where, on the plane of channel number against instrumental noise, it stops being honest. Nothing in the fluctuation-analysis literature addresses that. What the classical arm of this paper is should also be said plainly: least squares on the deterministic mean current with the unitary current held fixed. Fluctuation analysis is never benchmarked here, and no comparison against the variance-mean regression should be read into any result below.

The criterion this paper supplies is checkable. Under a correctly specified likelihood the score, the gradient of the log-likelihood with respect to the parameters, averages to zero at the true parameters, and the covariance of the score equals the Fisher information the likelihood reports, which is the information matrix equality \citep{white1982maximum,huber1967behavior}. Under an approximation both statements fail, and how they fail is measurable. Goodness of fit does not answer the question, because an approximation can reproduce the mean current closely and still report an information matrix wrong by more than an order of magnitude, and comparing two approximations to each other does not answer it either, because both may be wrong. What makes the question answerable for this class of models is an asymmetry between simulation and inference: the forward process can be simulated exactly while its likelihood cannot be computed exactly, so the simulator supplies the ground truth that the likelihood is an approximation of, and any candidate likelihood can be interrogated against data drawn from the very process it claims to describe. The sensitivity matrix is then analytic and the variability matrix is a Monte Carlo average over independent replicates, with no null hypothesis to reject and no finite-sample defect to correct for. This is measurement rather than testing.

The mismatch between the two matrices is itself the measurement. It gives an information distortion matrix that says, in units of parameter variance a reader already understands, by how much and in which parameter directions an approximation misreports what the data determine. The distortion runs in both directions: some approximations state more confidence than the data support, and others state less. % TODO-SIGN: direction convention as fixed in Diagnostics; if the pending code read reverses it, flip every TODO-SIGN site in one pass.
Evaluated across channel number and instrumental noise at a fixed acquisition interval, the same quantity yields a map of which parameters a macroscopic recording can deliver and whose reported interval is honest, with boundaries drawn as level sets of a continuous diagnostic rather than as thresholds. The map carries an ordering that is falsifiable and does not depend on where those level sets are placed: there is no region of the measured plane in which the amplitude still carries information and the classical error bar is simultaneously trustworthy, the two boundaries being separated by one to three decades of instrumental noise everywhere they were measured. % src: projects/eLife_2025/figures/paper_both/figure_6.html (crossing table and fitted slopes, Delta.k_off = 0.1)
Wherever the fluctuations still say something about the unitary current or the channel number, a classical fit reports an interval that is too narrow. % TODO-SIGN

Against that history, what is new here is narrow and is worth stating narrowly. A joint estimate of rates and amplitudes from fluctuation data, carrying standard errors, is \citet{anderson1973voltage}. What this paper adds, to a non-stationary protocol read from a single record, is the third element of the combination: an interval whose calibration has been verified rather than assumed. The thousands of replicate recordings used below certify the method; they are not what applying it requires.

The measurements are made in a deliberately small setting: a two-state channel whose open probability at the peak of the response is fixed at one half, a single concentration jump, simulated data, and comparison at the level of the likelihood rather than of a fitted experiment. The small model isolates the error the approximation itself commits, with no confounding from a misspecified kinetic scheme, and it keeps the exact simulation cheap enough to run thousands of replicates at every design cell. Richer schemes, the stationary regime, the few-channel boundary and experimental recordings are the subjects of companion work.

The map answers a question the field is currently arguing about. \citet{delcore2025parameter} argue against filtering on cost, since the moment integration and the correction step run between every pair of samples, and they are right about the cost. Knowing which approximation misreports the uncertainty, by how much, and in which corner of the operating regime is what turns that cost into a decision a reader can make for their own rig.
